% =========================================================
% Proof: Fraction Pair Equivalence Relation
% Source: notes/rational-numbers/notes-rational-numbers.tex
% Difficulty: Easy
% =========================================================

\subsection*{Fraction Pair Equivalence Relation}
\label{prf:fraction-pair-equivalence}

\begin{remark*}[Return]
$\leftarrow$ \hyperref[thm:fraction-pair-equivalence]{Back to Theorem in Notes}
\end{remark*}

\begin{proof}

\Claim The relation on
\[
W = \{(a,b) : a,b \in \mathbb{Z},\ b \neq 0\}
\]
defined by
\[
(a,b) \sim (c,d) \iff ad = bc
\]
is an equivalence relation. Consequently, its equivalence classes
partition $W$ and represent the rational numbers.

\Given $(a,b),(c,d),(e,f) \in W$.

\Goal Show that $\sim$ is reflexive, symmetric, and transitive.

\Strategy
\begin{itemize}
\item Reflexive: show $ab = ba$.
\item Symmetric: rearrange $ad = bc$.
\item Transitive: combine $ad = bc$ and $cf = de$ to obtain $af = be$.
\end{itemize}

% -------------------------------------------------------
% YOUR PROOF HERE
% -------------------------------------------------------

\end{proof}

\begin{remark*}[Proof shape]
Each part follows from elementary properties of integer arithmetic.
The transitive step requires combining two cross-product equalities.
\end{remark*}

\begin{remark*}[Dependencies]
\begin{itemize}
  \item Definition~\ref{def:fraction-pair-equality}: relation $ad = bc$.
  \item Definition of equivalence relation (reflexive, symmetric, transitive).
  \item Basic properties of integer multiplication.
\end{itemize}
\end{remark*}